\scanpage{86}
% archive sheet number: 85
% L1--2
$(Q\,\mathrm{Sch})$ \unsure{Catégorie obtenue à partir de} $(\mathrm{Sch})$ \unsure{en rendant inversibles les morphismes entiers radiciels surjectifs}

% L3--6
\unsure{Prétopologies} fppf et fpqc définies ; \unsure{on parlera de faisceaux|on parle de faisceaux} q.fpqc et q.fppf, \unsure{ce sont les} faisceaux fpqc et fppf qui \unsure{transforment morphismes entiers surj. rad. en} isom.

% Large category diagram. Thick headless strokes are kept headless.
\diagram{\begin{tikzcd}[ampersand replacement=\&,column sep=small,row sep=large,scale=.65,transform shape,
execute at begin picture={
\draw[line width=.8pt] (-6.4,4.2)--(1.4,4.2)--(1.4,3.2);
\draw (-6.4,3.7)--(-6.4,-.7)--(-4.1,-.7);
\draw (-3.8,-.7)--(-3.8,3.0)--(-1.8,3.0);
\draw (-1.8,3.0)--(-1.8,-3.5)--(1.0,-3.5)--(1.0,2.8)--(1.7,2.8);
}]
\&\&\begin{gathered}\text{faisceaux sur le}\text{ site des $k$-schémas}\\\text{qcts parfaits}\end{gathered}\arrow[dr,dash]\&\\
\begin{gathered}\text{groupes de la cat.}\text{ des $k$-schémas}\\\text{loc. qcts. parfaits}\end{gathered}\arrow[urr]\arrow[dr,dash]\arrow[dd,dash]
\&\begin{gathered}\text{\deleted{\unsure{groupes de la catégorie des}}}\\\text{\deleted{\unsure{$k$-schémas loc. qcts parfaits}}}\end{gathered}
\&\&\begin{gathered}\text{faisceaux sur le site}\\\text{\deleted{\unsure{q.fppf}} des $k$-schémas}\\\text{\deleted{loc.} qcts. parfaits}\\\text{essentiellement de}\text{ t.f. sur $k$}\end{gathered}\arrow[dd,dash]\\
\&\text{$k$-sch. en gr./rad.}\arrow[uur]\arrow[d,dash]\arrow[dr,dash]\&\&\\
\begin{gathered}\text{groupes de la catégorie}\text{ des $k$-schémas}\\\text{qcts parfaits}\end{gathered}\arrow[r,dash]
\&\begin{gathered}\text{$k$-gr. qp. qcts,}\\\text{mod rad}\end{gathered}\arrow[dr,dash]
\&\begin{gathered}\text{$k$-gr. loc. alg./radiciels}\\=\text{\deleted{\unsure{groupes algébriques}}}\\\text{\deleted{\unsure{localement}}}\end{gathered}\arrow[d,dash]\arrow[uur,dash]
\&\text{gr. proquasialg.}\arrow[d,dash]\\
\&\&\begin{gathered}\text{$k$-gr. alg./radiciels}\\\simeq\text{groupes quasi alg.}\\\text{ess de t.f.}\end{gathered}\arrow[r,dash]
\&\begin{gathered}\text{$k$-gr qp. qcts/radiciels}\\\text{\deleted{\unsure{mod rad}}}\end{gathered}
\end{tikzcd}}
